
\subsubsection{\textit{Physical Consequence: The Fundamental Mass Quantum (Electron Mass).}}
$Z_{MAR}$ establishes the theoretical resolution floor for a stable topological boundary. Because mass is fundamentally the resistance to state transition (inertia), this geometric noise floor manifests physically as the lowest stable massive excitation: the electron. Any charged configuration with a coupling lighter than this threshold cannot maintain stability against the vacuum flux. This fixes the mass \textit{ratio}:
\begin{equation}
\frac{m_e}{M_P} \propto \frac{Z_{MAR}}{Z_{geo}}
\end{equation}
The Planck/electron mass hierarchy is therefore a geometric consequence of the substrate's signal-to-noise ratio.

\subsubsection{\textit{Physical Consequence: The Scale of Matter (Atomic Length).}}
This resolution floor simultaneously determines the spatial scale of the periodic table. The Bohr Radius ($a_0$) emerges as the ratio of the vacuum's geometric impedance (signal strength) to the electron's mass resolution (noise floor).

The fine-structure constant is the inverse geometric impedance ($\alpha = 1/Z_{geo}$). Substituting these structural definitions into the Bohr radius formula reveals the vacuum's \textbf{Dynamic Range}, the signal-to-noise ratio of the lattice:
\begin{equation}
a_0 = \frac{\hbar}{c} \cdot \frac{Z_{geo}}{m_e} \propto \frac{Z_{geo}}{Z_{MAR}}
\end{equation}
(In natural units $\hbar = c = 1$, we have $a_0 = Z_{geo}/Z_{MAR}$.)

Because the electron mass is fixed by the persistence margin ($MAR$) and the coupling is fixed by the lattice topology, the \textbf{fundamental scale of chemistry} is structurally locked: the Angstrom scale ($10^{-10}$ m) is the immutable theater of atomic interaction. While complex atoms vary in effective radius, the underlying unit of atomic architecture is fixed by the lattice resolution.
